%\chapter*{Введение}
\chap{Введение}

Целью данной лабораторной работы является:
\begin{itemize}
	\item изучение применения древовидных структур данных для быстрого поиска и обработки запросов; 
	\item освоение операций с бинарным деревом поиска, понимание влияния высоты дерева на сложность алгоритмов и принцип балансировки; 
	\item получение практических навыков использования префиксных сумм, дерева Фенвика и дерева отрезков для эффективной обработки запросов к диапазонам данных.
\end{itemize}
%Целью данной лабораторной работы является изучение применения древовидных структур данных для быстрого поиска и обработки запросов; освоение операций с бинарным деревом поиска, понимание влияния высоты дерева на сложность алгоритмов и принцип балансировки; получение практических навыков использования префиксных сумм, дерева Фенвика и дерева отрезков для эффективной обработки запросов к диапазонам данных.
%\newline

Задачи:
\begin{enumerate}
	\item реализовать бинарное дерево поиска и основные операции с ним;
	\item исследовать влияние порядка вставки на высоту дерева;
	\item познакомиться с балансировкой на примере AVL-дерева;
	\item сравнить несколько подходов к вычислению диапазонных запросов и точечных обновлений.
\end{enumerate}

\chap{Выполнение}

\sect{Задание 1. Создание бинарного дерева}

\begin{figure}[!htbp]
	\centering
	%\includegraphics[width=0.6\linewidth]{""}
	\caption{}
\end{figure}
\FloatBarrier

\sect{Задание 2. Вставка в бинарное дерево поиска}

\sect{Задание 3. Поиск элемента}

\sect{Задание 4. Обходы дерева}

\sect{Задание 5. Удаление из BST}

\sect{Задание 6. Влияние порядка вставки на высоту дерева}

\sect{Задание 7. Повороты AVL-дерева}

\sect{Задание 8. Наивный запрос суммы диапазона}

\sect{Задание 9. Префиксные суммы}

\sect{Задание 10. Дерево Фенвика}

\sect{Задание 11. Дерево отрезков для суммы	}

\sect{Задание 12. Сумма, минимум и максимум}

\sect{Задание 13. Самостоятельная работа}

%\lstinputlisting[
%	style=python, 
%	caption={},
%	linerange={1-9}
%]{main.py}

\chap{Вывод}

В ходе лабораторной работы были:
\begin{itemize}
	\item изучены применения древовидных структур данных для быстрого поиска и обработки запросов; 
	\item освоены операции с бинарным деревом поиска, получено понимание влияния высоты дерева на сложность алгоритмов и принцип балансировки; 
	\item получены практические навыки использования префиксных сумм, дерева Фенвика и дерева отрезков для эффективной обработки запросов к диапазонам данных.
\end{itemize}

\chap{Контрольные вопросы}
\begin{enumerate}
	\item \textit{Что называется деревом и какие основные элементы используются при его описании?}\newline
	\item \textit{Что такое высота дерева и как она влияет на стоимость операций?}\newline
	\item \textit{Чем бинарное дерево поиска отличается от произвольного бинарного дерева?}\newline
	\item \textit{Как выполняются поиск и вставка ключа в BST?}\newline
	\item \textit{Какие случаи необходимо учитывать при удалении узла из BST?}\newline
	\item \textit{Чем отличаются preorder, inorder и postorder?}\newline
	\item \textit{Почему inorder-обход BST возвращает ключи в отсортированном порядке?}\newline
	\item \textit{Почему BST может выродиться при неудачном порядке вставки?}\newline
	\item \textit{Что такое баланс-фактор AVL-узла?}\newline
	\item \textit{Для чего используются одинарные и двойные повороты AVL-дерева?}\newline
	\item \textit{Чем в общем виде красно-чёрное дерево отличается от AVL-дерева?}\newline
	\item \textit{Почему B-Tree и B+Tree удобны для внешней памяти?}\newline
	\item \textit{В чём основная особенность организации данных в B+Tree?}\newline
	\item \textit{Что называется диапазонным запросом и зачем нужна предварительная обработка?}\newline
	\item \textit{Как строятся префиксные суммы и как с их помощью получить сумму диапазона?}\newline
	\item \textit{Почему префиксные суммы неудобны при частых изменениях массива?}\newline
	\item \textit{Для чего используется дерево Фенвика?}\newline
	\item \textit{Что определяет выражение i \& -i в дереве Фенвика?}\newline
	\item \textit{Как получить сумму диапазона с помощью дерева Фенвика?}\newline
	\item \textit{Что хранит узел дерева отрезков?}\newline
	\item \textit{Как выполняются диапазонный запрос и точечное обновление дерева отрезков?}\newline
	\item \textit{В каких случаях удобнее применять префиксные суммы, дерево Фенвика и дерево отрезков?}\newline
\end{enumerate}